%قبل از شروع مسئله ترجیحاً چیزی ننویسید. اما اگر خواستید بنویسید، در قالب لیستی (فقرات و شمارش) بنویسید.
%برای اضافه کردن نمره، بعدی از شروع‌مسئله، [x] را اضافه کنید و به جای x  نمره مورد نظر را قرار دهید. برای زیرمسئله‌ها هم به همین صورت است.
%\شروع‌مسئله[10]
\شروع‌مسئله[15]
\متن‌لاتین{RSA}
\پانویس{Rivest-Shamir-Adelman} یک سیستم رمزنگاری کلید عمومی\پانویس{public-key} است. این روش یکی از بزرگ‌ترین پیشرفت‌ها در زمینهٔ رمزنگاری و یکی از اساس‌های امنیت اینترنت به شمار می‌رود. \متن‌لاتین{RSA} روشی است که امنیت نشست‌های\پانویس{session} \متن‌لاتین{HTTPS}\پانویس{Hypertext Transfer Protocol Secure} سایت‌های تجارت الکترونیکی و خدمات ابری را فراهم می‌کند. این روش  در پروتکل‌های \متن‌لاتین{SSL}\پانویس{Secure Sockets Layer} و \متن‌لاتین{SSH}\پانویس{Secure Shell} نیز استفاده می‌شود.

در این سیستم به طور کلی هر کاربر دو عدد اول تصادفی بزرگ مانند $p$ و $q$ تولید می‌کند. این کاربر پیمانۀ عمومی\پانویس{public modulus} خود را برابر $m = p \cdot q$ می‌گیرد. پس از آن کاربر یک عدد کوچک به نام $e$ در نظر می‌گیرد که نسبت  به 
$(p-1)(q-1)$
اول باشد و $d$ را برابر با وارون ضربی $e$ نسبت به این پیمانه می‌گیرد. به عبارت دیگر 
$d\cdot e = 1 \mod (p-1)(q-1)$. این کاربر $(e,m)$ را به عنوان کلید عمومی خود اعلام می‌کند و $d$ را (به عنوان بخشی از کلید خصوصی) پیش خود نگه می‌دارد. اگر کسی بخواهد به این کاربر پیغامی مثل عدد $n$ ارسال کند، آن را به صورت 
$c = E(n) = n^e \mod m$
رمزگذاری می‌کند. کاربر نیز این پیغام را به صورت  $D(c) = c^d \mod m$ رمز گشایی می‌کند. بنابر قضیۀ اویلر تقریبا برای تمام $n$ها داریم 
$D(E(n)) = n$
که نشان می‌دهد پیام به روش گفته شده رمز گشایی می‌شود. با توجه به این‌که دیگر افراد عدد $d$ را ندارند، قادر به رمزگشایی پیام نخواهند بود (مگر این که به طریقی بتوانند $m$ را به عوامل اولش تجزیه کنند که راه سریعی برای این کار تاکنون پیدا نشده است).

\شروع{زیرمسائل}

\زیرمسئله[7]
\sloppy
فرض کنید $m$ یک عدد $k$ بیتی باشد و برای محاسبۀ $E(n)$ از الگوریتم \متن‌لاتین{MODULAR-EXPONENTIATION} که صفحۀ ۹۵۷ کتاب \متن‌لاتین{CLRS} (ویرایش سوم) آمده استفاده می‌کنیم. همچنین برای ضرب دو عدد از الگوریتم کاراتسوبا که در کلاس گفته شد استفاده می‌کنیم. در این صورت مرتبۀ زمانی رمزنگاری (محاسبۀ $E(n)$) را محاسبه کنید.

\پاسخ{
فرض کنید $d=\log_2 m$ باشد. الگوریتم \متن‌لاتین{MODULAR-EXPONENTIATION} از $O(d)$ ضرب استفاده می‌کند. اگر هر ضرب را با الگوریتم کاراتسوبا در $O(d^{\log_2 3})$ انجام دهیم، زمان اجرای رمزنگاری $O(d^{1+\log_2 3})$ است.

}


\زیرمسئله[8]
فرض کنید برای پیدا کردن اعداد اول ($p$ یا $q$) از الگوریتم تصادفی تست فرما
 (که در \href{https://en.wikipedia.org/wiki/Primality_test#Fermat_primality_test}{این لینک} توضیح داده شده)‌ برای بررسی اول بودن اعداد استفاده می‌کنیم. اگر از برخی پیچیدگی‌های الگوریتم صرف‌نظر کنیم، الگوریتم یک عدد تصادفی $k$ رقمی مثل $n$ انتخاب می‌شود و می‌خواهیم به صورت تصادفی بررسی کنیم که آیا این عدد اول است یا خیر. برای این بررسی مثلاً در ۱۰۰ مرحله یک عدد تصادفی $a$ انتخاب می‌شود و $a^{n-1} \mod n$ محاسبه می‌شود. اگر این عدد، برابر با یک نبود $n$ حتماً مرکب است. همچنین اگر $n$ یک عدد اول نباشد (با صرف‌نظر از اعداد کارمایکل)، برای نیمی از $a$های ممکن این عدد نامساوی با یک است.
\begin{itemize}
\فقره
این الگوریتم از نوع مونت‌کارلو است یا لاس‌وگاس؟ احتمال خطای آن را بیان کنید.
\فقره
با استفاده از قضیۀ عدد اول 
(که در \href{https://en.wikipedia.org/wiki/Prime_number_theorem}{این لینک} توضیح داده شده) احتمال این که عدد $n$ انتخاب شده اول باشد (به صورت مجانبی) چقدر است؟
\فقره
زمان اجرای میانگین این روش را برای پیدا کردن یک عدد اول $k$ رقمی بیان کنید.
\end{itemize}

\solution{
الگوریتم بررسی اول بودن یک عدد مونت‌کارلو است زیرا سریع است ولی احتمالاً درست است. توجه کنیم که این الگوریتم امکان ندارد عددی اول را مرکب اعلام کند، ولی این احتمال وجود دارد که عددی مرکب را اول اعلام کند (البته توجه کنیم که این الگوریتم به هیچ صورت نمی‌تواند اعداد کارمایکل را از اعداد اول تشخیص دهد و آن‌ها را در راه‌حل این سؤال در نظر نمی‌گیریم).
خطای این الگوریتم آن است که ممکن است عددی مرکب را اول در نظر بگیرد. بنا به صورت سؤال  می‌دانیم برای حداقل نیمی از 
$a$
‌های ممکن 
$a^{n-1} \neq 1 \mod n$.
پس (از آن‌جا که 
$100$
مرحله از یکدیگر مستقل هستند داریم)
\begin{equation*}
\text{احتمال خطا} = \text{احتمال خطا در یک مرحله}^{100} \leq (\frac{1}{2})^{100} = 2^{-100}
\end{equation*} 
 بنا بر قضیۀ عدد اول، اگر $pi(n)$ تعداد اعداد اول کوچکتر مساوی $n$ باشد داریم $\pi(n) \sim \frac{n}{\ln n}$. بنابراین احتمال اول بودن $n$ از مرتبه $\pi(n)/n = 1/\ln n$ است.
 
 برای پیدا کردن میانگین زمان اجرای الگوریتم دقت کنید احتمال آن که این الگوریتم یک عدد را انتخاب و آن را اول اعلام کند با احتساب خطا نیز $1/\ln n$ است.
در نتیجه امید ریاضی مراحل لازم برای یافتن عدد اول برابر است با
$\frac{1}{O(\frac{1}{\ln n})} = O(\ln n) = O(d)$.
همچنین توجه کنیم برای هر تست $100$ بار عملیات به توان رساندن را انجام می‌دهیم که با توجه به قسمت قبل به توان رساندن در زمان 
$\mathcal{O}(k ^ {1 + \log_2 3})$
انجام می‌شود و ۱۰۰ هم عددی ثابت است. پس کل زمان لازم به صورت مجانبی برابر است با
$\mathcal{O}(k ^ {2 + \log_2 3}).$

}

\پایان{زیرمسائل}





%%%%%%%%%%%%%%%%%%%%%%%%%%%%%%%%%%%%%%%%%%%%%%%%%%%%%
% برخی دستورات مفید برای تمرین‌ها. %
%%%%%%%%%%%%%%%%%%%%%%%%%%%%%%%%%%%%%%%%%%%%%%%%%%%%%

% فرمول‌های نمونه:
%$f(x) = \Theta(x)$
%$T(x, y) \leq \log(x) + 2^y + \binom{2n}{n}$
% {\tt code\_function}


% لیست موارد بدون شماره‌گذاری:
%\شروع{فقرات}
%    \فقره مورد اول 
%        \شروع{فقرات}
%            \فقره مورد اول داخلی 
%                \شروع{فقرات}
%                    \فقره اولین مورد خیلی داخلی 
%                    \فقره دومین مورد خیلی داخلی 
%                \پایان{فقرات}
%            \فقره مورد دوم داخلی 
%        \پایان{فقرات}
%    \فقره مورد دوم 
%\پایان{فقرات}

% لیست موارد شماره‌گذاری شده:
%\شروع{شمارش}
%    \فقره یک مورد 
%    \فقره یک مورد دیگر 
%    \فقره باز هم یک مورد دیگر
%\پایان{شمارش}

% برای نوشتن یک فرمول چند خطی که یک شماره بخوره. برای حالت‌های دیگر به سایت زیر مراجعه کنید:
% https://www.overleaf.com/learn/latex/Aligning_equations_with_amsmath
%\begin{align} 
%    \begin{split}
%        (x+y)^3 &= (x+y)^2(x+y) \\
%                &= (x^2+2xy+y^2)(x+y) \\
%                &= (x^3+2x^2y+xy^2) + (x^2y+2xy^2+y^3) \\
%                &= x^3+3x^2y+3xy^2+y^3
%    \end{split}                                 
%\end{align}

% نوشتن ماتریس و آرایه دوبعدی
%\begin{align}
%    A = 
%    \begin{bmatrix}
%        A_{11} & A_{21} \\
%        A_{21} & A_{22}
%    \end{bmatrix}
%\end{align}

% اضافه کردن تصویر با شرح. برای اضافه کردن تصویر بدون شرح، فقط دستور درج‌تصویر را استفاده کنید.
%\شروع{شکل}[h!]
%\تنظیم‌ازوسط
%  \درج‌تصویر[پهنا=0.5\پهنای‌متن]{img.jpg}
%\شرح{شرح تصویر}
%\برچسب{تصویر:۱}
%\پایان{شکل}
%\پایان{زیرمسائل}

%کد زیر را ببینید.
%%\lstinputlisting{source_filename.py}
%\begin{latin}
%\begin{lstlisting}
%# Import the modules
%import sys
%import random
%
%ans = True
%
%while ans:
%question = raw_input("Ask the magic 8 ball a question: (press enter to quit) ")
%
%answers = random.randint(1,8)
%
%if question == "":
%sys.exit()
%
%elif answers == 1:
%print "It is certain"
%
%elif answers == 2:
%print "Outlook good"
%
%elif answers == 3:
%print "You may rely on it"
%
%elif answers == 4:
%print "Ask again later"
%
%elif answers == 5:
%print "Concentrate and ask again"
%
%elif answers == 6:
%print "Reply hazy, try again"
%
%elif answers == 7:
%print "My reply is no"
%
%elif answers == 8:
%print "My sources say no"
%\end{lstlisting}
%\end{latin}
